% bismi ALLAH ARRAHMAN ARRAHIM





























































\documentclass[11pt,a4paper]{article}

% =========================================================
% EX9 — Architecture de calcul modulaire à cores de sondes
% Version révisée : coordination neuronale adaptative
% =========================================================

\usepackage[utf8]{inputenc}
\usepackage[T1]{fontenc}
\usepackage[french]{babel}
\usepackage{lmodern}
\usepackage{microtype}

\usepackage{geometry}
\geometry{
  left=2.3cm,
  right=2.3cm,
  top=2.2cm,
  bottom=2.2cm
}

\usepackage{amsmath,amssymb,mathtools}
\usepackage{booktabs}
\usepackage{array}
\usepackage{longtable}
\usepackage{tabularx}
\usepackage{enumitem}
\usepackage{xcolor}
\usepackage{graphicx}
\usepackage{float}
\usepackage{caption}
\usepackage{hyperref}

\usepackage{tikz}
\usetikzlibrary{
  arrows.meta,
  positioning,
  fit,
  calc,
  shapes.geometric,
  shapes.misc,
  backgrounds,
  matrix
}

\usepackage{pgfplots}
\pgfplotsset{compat=1.18}

\hypersetup{
  colorlinks=true,
  linkcolor=blue!50!black,
  citecolor=blue!50!black,
  urlcolor=blue!50!black,
  pdftitle={EX9 -- Architecture modulaire à cores de sondes},
  pdfauthor={Projet EX}
}

\setlength{\parindent}{0pt}
\setlength{\parskip}{5pt}
\setlist{nosep}

% ---------------------------------------------------------
% Styles TikZ
% ---------------------------------------------------------

\tikzset{
  probe/.style={
    draw=blue!65!black,
    fill=blue!7,
    rounded corners=2pt,
    minimum width=1.75cm,
    minimum height=0.72cm,
    align=center,
    font=\small
  },
  core/.style={
    draw=black!75,
    fill=gray!8,
    rounded corners=5pt,
    thick,
    inner sep=8pt
  },
  neuralcore/.style={
    draw=orange!75!black,
    fill=orange!8,
    rounded corners=5pt,
    thick,
    inner sep=8pt
  },
  control/.style={
    draw=orange!70!black,
    fill=orange!12,
    rounded corners=3pt,
    minimum width=2.7cm,
    minimum height=0.8cm,
    align=center,
    font=\small
  },
  memory/.style={
    draw=green!50!black,
    fill=green!8,
    cylinder,
    shape border rotate=90,
    aspect=0.25,
    minimum width=1.6cm,
    minimum height=1.0cm,
    align=center,
    font=\small
  },
  data/.style={
    draw=purple!65!black,
    fill=purple!8,
    rounded corners=3pt,
    minimum width=1.8cm,
    minimum height=0.7cm,
    align=center,
    font=\small
  },
  flow/.style={
    -{Stealth[length=2mm]},
    thick
  },
  comm/.style={
    {Stealth[length=2mm]}-{Stealth[length=2mm]},
    dashed,
    thick
  },
  learned/.style={
    -{Stealth[length=2mm]},
    thick,
    orange!80!black
  },
  dependency/.style={
    -{Stealth[length=2mm]},
    thick,
    blue!70!black
  }
}

\title{
  \textbf{EX9}\\[0.3cm]
  \Large Architecture de calcul modulaire\\
  \Large à cores composés de sondes spécialisées\\[0.2cm]
  \normalsize
  Coordination neuronale adaptative et graphes de calcul dynamiques
}

\author{Projet EX}
\date{Document de conception et plan expérimental}

\begin{document}

\maketitle

\begin{abstract}

EX9 propose une architecture de calcul modulaire dans laquelle les
unités fonctionnelles sont organisées en cores spécialisés, chaque
core étant lui-même composé de sondes.

L'architecture distingue trois formes de réalisation :
(i) les cores logiques, organisés dans un graphe logiciel ;
(ii) les cores physiques, associés à des ressources matérielles ;
(iii) les cores neuronaux, construits à partir de sondes adaptatives
et destinés notamment à remplacer des blocs MLP.

L'architecture comprend plusieurs fonctions spécialisées :
calcul, composition, coordination, stockage temporaire et
communication.

Une caractéristique centrale d'EX9 est l'introduction d'un
\textbf{coordinateur neuronal adaptatif}. Celui-ci n'est pas un
contrôleur externe à règles fixes. Il constitue lui-même un core
composé de sondes adaptatives, dont les paramètres sont appris.

Ce coordinateur observe l'état du système, les données disponibles
et le graphe de calcul. Il apprend à sélectionner les cores,
ordonner les opérations, organiser les échanges, réutiliser les
résultats et, lorsque cela est autorisé, modifier dynamiquement
la structure du graphe d'exécution.

L'hypothèse scientifique est qu'une organisation de ce type pourrait
adapter la structure du calcul à chaque entrée ou situation, plutôt
que d'imposer systématiquement le même chemin d'exécution.

Les gains éventuels en précision, latence, mémoire et énergie ne sont
pas considérés comme acquis. Ils devront être établis par des
expériences contrôlées comparant EX9 à des architectures de référence.

\end{abstract}

\tableofcontents
\newpage

% =========================================================
\section{Positionnement et continuité avec EX7}
% =========================================================

\subsection{Point de départ expérimental}

EX7 V003, désigné \textit{OuterCoord-Lite-7x1}, constitue le point
de départ architectural disponible.

Cette version comprend :
\begin{itemize}
  \item sept sondes indépendantes de type \texttt{AdaptiveRBF} ;
  \item un coordinateur global ;
  \item une pondération globale apprise, notée $\rho$ ;
  \item une entrée commune transmise aux sept sondes ;
  \item une fusion des sorties des sondes.
\end{itemize}

Chaque sonde possède notamment :
\begin{itemize}
  \item un encodeur appris ;
  \item des centres RBF ;
  \item des largeurs adaptatives ;
  \item une couche de lecture produisant les logits.
\end{itemize}

Dans EX7 V003, les sept sondes sont entraînées sur des partitions
distinctes des exemples d'entraînement. Elles reçoivent néanmoins
toutes les mêmes dimensions d'entrée à l'inférence.

Il est donc important de distinguer deux formes de partitionnement :

\begin{itemize}
  \item le partitionnement des exemples d'entraînement entre les sondes ;
  \item le partitionnement éventuel des dimensions d'entrée, qui
        n'est pas réalisé dans EX7 V003.
\end{itemize}

Les sept sondes sont exécutées à chaque inférence. La pondération
$\rho$ est une pondération globale apprise, mais elle ne constitue
pas une politique de coordination dépendante de chaque entrée.

EX7 démontre ainsi la faisabilité d'un ensemble de sondes adaptatives
coordonnées. Il ne démontre pas encore l'intérêt d'un coordinateur
neuronal dynamique, d'une organisation hiérarchique en cores ou
d'une modification du graphe de calcul.

\subsection{Évolution architecturale proposée}

EX9 introduit une organisation hiérarchique :

\[
\boxed{
\text{Sondes}
\longrightarrow
\text{Cores spécialisés}
\longrightarrow
\text{Coordinateur neuronal}
\longrightarrow
\text{Système EX9}
}
\]

Un core est un ensemble structuré de sondes qui coopèrent pour
réaliser une fonction identifiable.

Le coordinateur est lui-même un core. Il est composé de sondes
adaptatives et apprend à organiser les autres cores.

\subsection{Objectifs}

Les objectifs d'EX9 sont les suivants :

\begin{enumerate}
  \item Décomposer le calcul en opérations spécialisées.
  \item Organiser les opérations sous forme de graphes composables.
  \item Regrouper les sondes en cores fonctionnels.
  \item Apprendre la coordination entre les cores.
  \item Permettre la sélection dynamique des opérations utiles.
  \item Apprendre à réutiliser les résultats intermédiaires.
  \item Étudier la modification dynamique du graphe de calcul.
  \item Construire des cores neuronaux capables de remplacer des
        blocs MLP.
  \item Étudier les correspondances entre organisation logique,
        neuronale et physique.
\end{enumerate}

% =========================================================
\section{Principes architecturaux}
% =========================================================

\subsection{Les éléments fondamentaux}

EX9 repose sur cinq éléments fonctionnels :

\begin{enumerate}
  \item \textbf{Sondes :} unités élémentaires de transformation,
        de calcul, de représentation ou de décision.
  \item \textbf{Cores :} ensembles organisés de sondes spécialisées.
  \item \textbf{Graphe :} représentation des dépendances et des
        possibilités de composition.
  \item \textbf{Mémoire :} stockage des états et résultats intermédiaires.
  \item \textbf{Coordinateur neuronal :} core adaptatif qui apprend
        à organiser l'exécution du système.
\end{enumerate}

\subsection{Séparation entre fonction et réalisation}

Il faut distinguer deux dimensions.

La première dimension décrit la fonction d'un core :
\begin{itemize}
  \item calcul ;
  \item composition ;
  \item coordination ;
  \item mémoire ;
  \item communication ;
  \item traitement neuronal.
\end{itemize}

La seconde décrit sa réalisation :
\begin{itemize}
  \item logique ;
  \item neuronale ;
  \item physique.
\end{itemize}

Ces dimensions ne sont pas exclusives.

Par exemple, le core de coordination peut être :
\begin{itemize}
  \item défini comme un graphe logiciel ;
  \item implémenté par des sondes neuronales ;
  \item exécuté sur un processeur ou une unité matérielle dédiée.
\end{itemize}

Un core neuronal peut donc être exécuté dans un core logique,
puis éventuellement être associé à des ressources physiques.

% =========================================================
\section{Définition formelle d'une sonde}
% =========================================================

\subsection{Sonde comme opérateur}

Soit un espace d'entrée $\mathcal{X}$ et un espace de sortie
$\mathcal{Y}$.

Une sonde est une application paramétrée :

\[
P_{\theta} :
\mathcal{X}
\rightarrow
\mathcal{Y}.
\]

Elle transforme une entrée en sortie selon ses paramètres
et sa fonction spécialisée.

Une sonde peut représenter :
\begin{itemize}
  \item une transformation numérique ;
  \item une transformation neuronale ;
  \item une opération de composition ;
  \item une décision de routage ;
  \item une opération de mémoire ;
  \item une décision de coordination.
\end{itemize}

Dans le cas d'une sonde spécialisée dans une opération numérique :

\[
P_{\theta}^{(k)}(x)
\simeq
\phi_k(x),
\]

où $\phi_k$ représente l'opération visée.

Exemples :

\[
\begin{aligned}
P^{(\mathrm{PROD})}(a,b) &\simeq ab,\\
P^{(\mathrm{EXP})}(c) &\simeq e^c,\\
P^{(\mathrm{DIV})}(u,v) &\simeq \frac{u}{v},\\
P^{(\mathrm{LOG})}(j) &\simeq \log(j).
\end{aligned}
\]

Le symbole $\simeq$ désigne une approximation éventuelle.

Une opération arithmétique peut aussi être réalisée par une primitive
exacte. Il n'est pas nécessaire que toutes les opérations soient
des approximations neuronales.

\subsection{Sonde adaptative de type RBF}

Une sonde neuronale peut utiliser une représentation latente :

\[
z = \tanh(Wx+b),
\]

où :

\[
W\in\mathbb{R}^{d\times n},
\qquad
b\in\mathbb{R}^{d}.
\]

Pour $K$ centres $c_k$, les activations radiales sont :

\[
r_k(z)
=
\exp\left(
-\frac{\|z-c_k\|_2^2}{2\sigma_k^2}
\right).
\]

La sortie peut être définie par :

\[
P_{\theta}(x)
=
R(x)V+\beta,
\]

avec :

\[
R(x)=[r_1(z),\ldots,r_K(z)].
\]

Les paramètres appris comprennent notamment :

\[
\theta=(W,b,C,\sigma,V,\beta).
\]

Cette formulation reprend le principe de la sonde
\texttt{AdaptiveRBF} utilisée dans EX7.

Elle ne signifie pas que toutes les sondes EX9 devront être des RBF.
Le choix de la famille de sondes dépend de la fonction à réaliser.

\subsection{Sonde de décision}

Une sonde peut également produire une représentation utilisée
par un coordinateur.

Par exemple :

\[
P_{\theta}^{(\mathrm{decision})}(s_t)
=
q_t,
\]

où $q_t$ représente un vecteur de scores ou de caractéristiques
relatives à l'état courant du système.

Une telle sonde ne réalise pas nécessairement le calcul final.
Elle contribue à la décision collective du core auquel elle appartient.

% =========================================================
\section{Définition formelle d'un core}
% =========================================================

\subsection{Définition générale}

Un core $C_i$ est un ensemble structuré de $m_i$ sondes :

\[
C_i =
\left(
\mathcal{P}_i,
\Gamma_i,
\mathcal{M}_i,
\Theta_i
\right),
\]

où :

\begin{itemize}
  \item $\mathcal{P}_i$ est l'ensemble des sondes ;
  \item $\Gamma_i$ décrit leur organisation et leurs dépendances ;
  \item $\mathcal{M}_i$ décrit les mécanismes de mémoire ou d'état ;
  \item $\Theta_i$ regroupe les paramètres du core.
\end{itemize}

Le core implémente une fonction :

\[
C_i :
\mathcal{X}_i
\times
\mathcal{S}_i
\rightarrow
\mathcal{Y}_i
\times
\mathcal{S}_i',
\]

où $\mathcal{S}_i$ représente éventuellement l'état interne.

Cette formulation permet de représenter aussi bien un core sans état
qu'un core possédant une mémoire temporaire.

\subsection{Structure interne}

\begin{figure}[H]
\centering
\begin{tikzpicture}[node distance=0.45cm]
  \node[probe] (p1) {Sonde 1};
  \node[probe,right=of p1] (p2) {Sonde 2};
  \node[probe,right=of p2] (p3) {Sonde 3};
  \node[probe,right=of p3] (p4) {Sonde $m$};

  \node[core,fit=(p1)(p2)(p3)(p4),
        label=above:{\textbf{CORE}}] {};

  \node[data,below=0.9cm of p2] (in) {Entrées};
  \node[data,right=5.1cm of in] (out) {Sorties};

  \draw[flow] (in.north) -- (p2.south);
  \draw[flow] (p3.south) -- (out.north);

  \draw[comm] (p1) -- (p2);
  \draw[comm] (p2) -- (p3);
  \draw[comm] (p3) -- (p4);
\end{tikzpicture}
\caption{Un core est composé de sondes qui coopèrent.}
\end{figure}

Cette représentation est conceptuelle. Les sondes ne sont pas
obligatoirement exécutées en série. Leur organisation peut être
parallèle, hiérarchique ou conditionnelle.

\subsection{Contrat fonctionnel d'un core}

Chaque core doit définir un contrat explicite :

\[
\mathcal{K}_i =
(\mathcal{X}_i,\mathcal{Y}_i,\mathcal{S}_i,
\mathcal{R}_i,\mathcal{C}_i),
\]

où :

\begin{itemize}
  \item $\mathcal{X}_i$ définit les entrées acceptées ;
  \item $\mathcal{Y}_i$ définit les sorties produites ;
  \item $\mathcal{S}_i$ définit l'état éventuellement requis ;
  \item $\mathcal{R}_i$ définit les ressources nécessaires ;
  \item $\mathcal{C}_i$ définit les contraintes d'exécution.
\end{itemize}

Le contrat permet au coordinateur de raisonner sur les cores
sans devoir connaître tous leurs détails internes.

% =========================================================
\section{Les trois formes de réalisation des cores}
% =========================================================

\subsection{Core logique}

Un core logique est une unité fonctionnelle logicielle constituée
de sondes et de mécanismes de composition.

Il n'impose aucune correspondance directe avec un coeur matériel.
Plusieurs cores logiques peuvent être exécutés sur un même CPU ou GPU.

Formellement :

\[
C_i^{(L)}
=
(\mathcal{P}_i,\Gamma_i,\mathcal{S}_i).
\]

Son graphe interne est :

\[
G_i=(V_i,E_i),
\]

où :
\begin{itemize}
  \item $V_i$ représente les sondes et opérations ;
  \item $E_i$ représente les dépendances de données et de contrôle.
\end{itemize}

\begin{figure}[H]
\centering
\begin{tikzpicture}[node distance=0.65cm]
  \node[data] (x) {Entrée};
  \node[probe,right=of x] (p1) {Sonde A};
  \node[probe,above right=0.5cm and 0.8cm of p1] (p2) {Sonde B};
  \node[probe,below right=0.5cm and 0.8cm of p1] (p3) {Sonde C};
  \node[probe,right=1.0cm of p2] (p4) {Sonde D};
  \node[data,right=1.0cm of p4] (y) {Sortie};

  \draw[flow] (x) -- (p1);
  \draw[flow] (p1) -- (p2);
  \draw[flow] (p1) -- (p3);
  \draw[flow] (p2) -- (p4);
  \draw[flow] (p3) -- (p4);
  \draw[flow] (p4) -- (y);

  \begin{scope}[on background layer]
    \node[core,fit=(p1)(p2)(p3)(p4),
          label=above:{Core logique}] {};
  \end{scope}
\end{tikzpicture}
\caption{Exemple d'un core logique organisé en graphe logiciel.}
\end{figure}

Le core logique définit l'organisation fonctionnelle, mais ne garantit
pas à lui seul une accélération matérielle.

\subsection{Core physique}

Un core physique correspond à une réalisation matérielle d'un ensemble
de sondes.

Il peut comprendre :
\begin{itemize}
  \item des unités de calcul ;
  \item des unités de composition ;
  \item une mémoire locale ;
  \item des mécanismes de communication ;
  \item un contrôleur d'exécution.
\end{itemize}

On définit :

\[
C_i^{(P)}
=
(\mathcal{H}_i,\mathcal{P}_i,\mathcal{M}_i,\mathcal{N}_i),
\]

où :
\begin{itemize}
  \item $\mathcal{H}_i$ désigne les ressources matérielles ;
  \item $\mathcal{P}_i$ désigne les sondes exécutées ;
  \item $\mathcal{M}_i$ désigne la mémoire locale ;
  \item $\mathcal{N}_i$ désigne le réseau de communication.
\end{itemize}

\begin{figure}[H]
\centering
\begin{tikzpicture}[node distance=0.5cm]
  \node[probe] (p1) {Unité sonde 1};
  \node[probe,right=of p1] (p2) {Unité sonde 2};
  \node[probe,right=of p2] (p3) {Unité sonde 3};

  \node[control,below=0.9cm of p2] (ctrl) {Contrôleur};
  \node[memory,below=0.8cm of p1] (mem) {Mémoire locale};
  \node[data,below=0.8cm of p3] (net) {Interface réseau};

  \node[core,fit=(p1)(p2)(p3)(ctrl)(mem)(net),
        label=above:{\textbf{CORE PHYSIQUE}}] {};

  \draw[comm] (p1) -- (p2);
  \draw[comm] (p2) -- (p3);
  \draw[flow] (ctrl) -- (p2);
  \draw[comm] (mem) -- (p1);
  \draw[comm] (net) -- (p3);
\end{tikzpicture}
\caption{Organisation conceptuelle d'un core physique.}
\end{figure}

Le passage à cette variante nécessite une étude de :
\begin{itemize}
  \item la surface matérielle ;
  \item la consommation énergétique ;
  \item la bande passante mémoire ;
  \item la latence des communications ;
  \item le parallélisme réellement exploitable.
\end{itemize}

Un core physique n'est pas nécessairement plus rapide qu'un CPU ou GPU.
Le bénéfice doit être mesuré sur une réalisation concrète.

\subsection{Core neuronal}

Un core neuronal est un ensemble de sondes adaptatives qui réalise
une transformation apprise.

Il est conçu notamment pour remplacer un bloc MLP ou une partie
d'un réseau neuronal classique.

On définit :

\[
C_i^{(N)}(x)
=
\mathcal{F}_{\Theta_i}
\left(
P_{i,1}(x),\ldots,P_{i,m_i}(x)
\right),
\]

où $\mathcal{F}_{\Theta_i}$ est l'opérateur de composition neuronal.

La fonction du core est apprise à partir de données :

\[
\Theta_i^\star
=
\arg\min_{\Theta_i}
\mathbb{E}_{(x,y)\sim\mathcal{D}}
\left[
\mathcal{L}
\left(
C_i^{(N)}(x),y
\right)
\right].
\]

\begin{figure}[H]
\centering
\begin{tikzpicture}[node distance=0.55cm]
  \node[data] (x) {Vecteur $x$};

  \node[probe,right=of x,yshift=0.9cm] (p1) {Sonde adaptative 1};
  \node[probe,right=of x] (p2) {Sonde adaptative 2};
  \node[probe,right=of x,yshift=-0.9cm] (p3) {Sonde adaptative 3};

  \node[control,right=1.0cm of p2] (fusion) {Composition\\neurale};
  \node[data,right=of fusion] (y) {Sortie};

  \draw[flow] (x) -- (p1);
  \draw[flow] (x) -- (p2);
  \draw[flow] (x) -- (p3);
  \draw[flow] (p1) -- (fusion);
  \draw[flow] (p2) -- (fusion);
  \draw[flow] (p3) -- (fusion);
  \draw[flow] (fusion) -- (y);

  \begin{scope}[on background layer]
    \node[core,fit=(p1)(p2)(p3)(fusion),
          label=above:{\textbf{CORE NEURONAL}}] {};
  \end{scope}
\end{tikzpicture}
\caption{Core neuronal composé de sondes adaptatives.}
\end{figure}

Une sonde n'est pas nécessairement un neurone scalaire. Une sonde
adaptative peut être un module vectoriel possédant son propre encodeur,
ses paramètres et sa fonction de sortie.

\subsection{Comparaison des trois formes}

\begin{table}[H]
\centering
\small
\begin{tabularx}{\textwidth}{p{2.5cm}XXX}
\toprule
 & Core logique & Core physique & Core neuronal \\
\midrule
Nature
& Logicielle
& Matérielle
& Apprise \\

Composition
& Sondes et graphe
& Sondes et unités matérielles
& Sondes adaptatives \\

Exécution
& CPU/GPU
& Unités physiques dédiées
& Calcul neuronal \\

Objectif
& Organiser
& Réaliser matériellement
& Apprendre une transformation \\

Validation
& Coût logiciel
& Latence, énergie, surface
& Qualité, paramètres, latence \\
\bottomrule
\end{tabularx}
\caption{Différences entre les formes de réalisation des cores.}
\end{table}

Ces formes sont complémentaires. Un core neuronal peut être défini
logiquement et exécuté physiquement.

% =========================================================
\section{Calcul sous forme de graphe de sondes}
% =========================================================

\subsection{Calcul classique}

Considérons :

\[
f(a,b,c,j)
=
ab+\frac{e^c}{a}+\log(j).
\]

Dans une expression classique, le calcul est décrit par une suite
d'opérations arithmétiques.

EX9 propose de représenter chaque opération par une sonde ou une
unité spécialisée.

\subsection{Graphe de calcul}

Définissons un graphe orienté :

\[
G=(V,E).
\]

Chaque sommet $v_i\in V$ représente une opération ou une sonde.
Chaque arête $(v_i,v_j)\in E$ indique qu'une sortie de $v_i$ est
utilisée par $v_j$.

Pour l'expression précédente :

\[
\begin{aligned}
u_1 &= P_{\mathrm{PROD}}(a,b),\\
u_2 &= P_{\mathrm{EXP}}(c),\\
u_3 &= P_{\mathrm{DIV}}(u_2,a),\\
u_4 &= P_{\mathrm{LOG}}(j),\\
y &= P_{\mathrm{SUM}}(u_1,u_3,u_4).
\end{aligned}
\]

Si les sondes sont exactes, cette composition réalise l'expression
arithmétique initiale. Si elles sont apprises, une erreur
d'approximation peut apparaître à chaque étape.

\begin{figure}[H]
\centering
\begin{tikzpicture}[node distance=0.75cm]
  \node[data] (ab) {$a,b$};
  \node[data,below=0.8cm of ab] (c) {$c$};
  \node[data,below=0.8cm of c] (j) {$j$};

  \node[probe,right=of ab] (prod) {PROD};
  \node[probe,right=of c] (exp) {EXP};
  \node[probe,right=of j] (log) {LOG};

  \node[probe,right=1.3cm of exp] (div) {DIV};
  \node[probe,right=1.3cm of div] (sum) {SUM};
  \node[data,right=of sum] (y) {$y$};

  \draw[flow] (ab) -- (prod);
  \draw[flow] (c) -- (exp);
  \draw[flow] (exp) -- (div);
  \draw[flow] (ab.east) -- ++(0.45,0) |- (div.north);
  \draw[flow] (j) -- (log);
  \draw[flow] (prod) -- (sum);
  \draw[flow] (div) -- (sum);
  \draw[flow] (log.east) -- ++(0.45,0) |- (sum.south);
  \draw[flow] (sum) -- (y);
\end{tikzpicture}
\caption{Graphe de calcul composé de sondes spécialisées.}
\end{figure}

\subsection{Dépendances et parallélisme}

Les opérations PROD, EXP et LOG sont indépendantes au départ.
Elles peuvent donc être exécutées en parallèle si les ressources
disponibles le permettent.

L'opération DIV dépend de EXP et de la valeur $a$.
L'opération SUM dépend des trois branches.

On note $\operatorname{Pred}(v)$ l'ensemble des prédécesseurs
d'un sommet $v$.

Une opération est exécutable lorsque toutes ses entrées sont
disponibles :

\[
\operatorname{Ready}(v,t)
=
\bigwedge_{u\in\operatorname{Pred}(v)}
\operatorname{Available}(u,t).
\]

Le graphe représente donc les dépendances réelles du calcul.

% =========================================================
\section{Composition des sondes et des cores}
% =========================================================

\subsection{Composition fonctionnelle}

Soient deux sondes :

\[
P_1:\mathcal{X}\rightarrow\mathcal{Z},
\qquad
P_2:\mathcal{Z}\rightarrow\mathcal{Y}.
\]

Leur composition est :

\[
P_{2\circ1}(x)=P_2(P_1(x)).
\]

Pour un graphe général, chaque sonde reçoit les sorties de ses
prédécesseurs :

\[
h_i
=
P_i
\left(
\{h_j:j\in\operatorname{Pred}(i)\}
\right).
\]

La sortie globale est produite par un ou plusieurs sommets terminaux.

\subsection{Composition hiérarchique}

Les sondes peuvent être regroupées en cores, puis les cores
peuvent être composés :

\[
\mathcal{P}
\longrightarrow
\{C_1,\ldots,C_M\}
\longrightarrow
F_{\mathrm{EX9}}.
\]

La fonction globale devient :

\[
F_{\mathrm{EX9}}(x)
=
\mathcal{C}
\left(
C_1(x_1),\ldots,C_M(x_M)
\right).
\]

Les entrées $x_i$ peuvent être identiques, projetées, partitionnées
ou partagées selon la topologie retenue.

La politique de coordination détermine quels cores sont activés
et comment leurs résultats sont composés.

% =========================================================
\section{Mémoire temporaire et réutilisation}
% =========================================================

\subsection{Problème de redondance}

Un même sous-calcul peut être utilisé plusieurs fois.

Par exemple :

\[
u=\exp(c)
\]

peut alimenter plusieurs opérations.

Une architecture sans mécanisme de réutilisation peut recalculer
plusieurs fois la même valeur.

EX9 propose un stockage temporaire des résultats intermédiaires.

\subsection{Mémoire temporaire}

On définit une mémoire :

\[
\mathcal{M}_t
=
\{(k,v,\tau,\ell)\},
\]

où :
\begin{itemize}
  \item $k$ est une clé identifiant le calcul ;
  \item $v$ est le résultat calculé ;
  \item $\tau$ est l'instant ou la version de calcul ;
  \item $\ell$ est la durée de validité ou la politique de conservation.
\end{itemize}

Une opération peut consulter la mémoire avant de lancer un calcul :

\[
\operatorname{Execute}(k,x)
=
\begin{cases}
\mathcal{M}_t[k],
& \text{si une valeur valide existe},\\
P_k(x),
& \text{sinon}.
\end{cases}
\]

Après calcul, le résultat peut être enregistré.

\subsection{Exemple de réutilisation}

\[
\begin{aligned}
u &= \operatorname{EXP}(c),\\
v_1 &= \operatorname{DIV}(u,a),\\
v_2 &= \operatorname{PROD}(u,b),\\
v_3 &= \operatorname{SUM}(v_1,v_2).
\end{aligned}
\]

La valeur $u$ est calculée une fois et réutilisée deux fois.

\begin{figure}[H]
\centering
\begin{tikzpicture}[node distance=0.8cm]
  \node[data] (c) {$c$};
  \node[probe,right=of c] (exp) {EXP};
  \node[memory,right=of exp] (mem) {Mémoire\\temporaire};

  \node[probe,right=1.5cm of mem,yshift=0.8cm] (div) {DIV};
  \node[probe,right=1.5cm of mem,yshift=-0.8cm] (prod) {PROD};

  \node[probe,right=1.2cm of div,yshift=-0.8cm] (sum) {SUM};

  \draw[flow] (c) -- (exp);
  \draw[flow] (exp) -- (mem);
  \draw[comm] (mem) -- (div);
  \draw[comm] (mem) -- (prod);
  \draw[flow] (div) -- (sum);
  \draw[flow] (prod) -- (sum);
\end{tikzpicture}
\caption{Une valeur intermédiaire est calculée une fois puis réutilisée.}
\end{figure}

\subsection{Conditions de validité}

La réutilisation n'est correcte que si les entrées et le contexte
du calcul sont compatibles.

La clé de mémoire doit prendre en compte, selon le cas :
\begin{itemize}
  \item les valeurs d'entrée ;
  \item les paramètres de la sonde ;
  \item la version du modèle ;
  \item le contexte numérique ;
  \item les éventuels effets de bord.
\end{itemize}

Pour des sondes stochastiques ou dépendantes d'un état mutable,
une réutilisation naïve peut modifier le comportement du système.

La mémoire doit donc être accompagnée d'une politique explicite
de validité et d'invalidation.

% =========================================================
\section{Le coordinateur neuronal adaptatif}
% =========================================================

\subsection{Principe central}

Dans EX9, le coordinateur n'est pas conçu comme un ensemble de règles
fixes écrites manuellement.

Il constitue un core neuronal à part entière, composé de sondes
adaptatives.

Il possède :
\begin{itemize}
  \item ses propres sondes ;
  \item ses propres paramètres appris ;
  \item une représentation de l'état du système ;
  \item une politique de décision ;
  \item un mécanisme de sélection d'actions ;
  \item un objectif d'apprentissage intégrant la qualité et le coût
        de l'exécution.
\end{itemize}

Le coordinateur apprend à organiser les autres cores.

Il peut notamment apprendre :
\begin{itemize}
  \item quels cores activer ;
  \item quelles opérations exécuter ;
  \item dans quel ordre les exécuter ;
  \item quelles opérations exécuter en parallèle ;
  \item quels résultats réutiliser ;
  \item quand stocker ou libérer un résultat ;
  \item comment router les informations ;
  \item quand modifier le graphe de calcul.
\end{itemize}

Le coordinateur ne se contente donc pas de choisir une sonde.
Il apprend une stratégie d'organisation du calcul.

\subsection{Définition du core de coordination}

On définit le core de coordination :

\[
C_{\mathrm{coord}}
=
\left(
\mathcal{P}_{\mathrm{coord}},
\Gamma_{\mathrm{coord}},
\mathcal{M}_{\mathrm{coord}},
\Theta_{\mathrm{coord}}
\right).
\]

Il est constitué de sondes adaptatives :

\[
\mathcal{P}_{\mathrm{coord}}
=
\{
P_1^{(c)},P_2^{(c)},\ldots,P_m^{(c)}
\}.
\]

Chaque sonde peut apprendre une spécialisation dans une fonction
de coordination, par exemple :
\begin{itemize}
  \item évaluer les cores disponibles ;
  \item estimer le coût d'une opération ;
  \item choisir une stratégie de composition ;
  \item décider d'une réutilisation ;
  \item proposer une modification du graphe ;
  \item évaluer les ressources nécessaires.
\end{itemize}

Ces spécialisations sont des possibilités architecturales.
Elles ne supposent pas que chaque fonction doive obligatoirement
être attribuée à une sonde unique.

\subsection{Entrées du coordinateur}

À l'instant $t$, le coordinateur reçoit un état :

\[
s_t =
(x_t,G_t,\mathcal{M}_t,\mathcal{R}_t,\mathcal{H}_t),
\]

où :

\begin{itemize}
  \item $x_t$ représente les données disponibles ;
  \item $G_t$ représente le graphe courant ;
  \item $\mathcal{M}_t$ représente la mémoire temporaire ;
  \item $\mathcal{R}_t$ représente les ressources disponibles ;
  \item $\mathcal{H}_t$ représente l'historique d'exécution utile.
\end{itemize}

Le coordinateur peut recevoir des représentations compactes de ces
éléments plutôt que leur contenu intégral.

Sa décision est :

\[
a_t =
C_{\mathrm{coord}}
(s_t;\Theta_{\mathrm{coord}}).
\]

Le coordinateur produit une action ou un ensemble d'actions.

\subsection{Politique neuronale}

On peut formaliser le coordinateur par une politique :

\[
\pi_{\theta_c}(a_t\mid s_t,G_t),
\]

où $\theta_c$ représente ses paramètres appris.

Dans une formulation probabiliste :

\[
a_t
\sim
\pi_{\theta_c}(\cdot\mid s_t,G_t).
\]

Dans une formulation déterministe :

\[
a_t
=
\pi_{\theta_c}(s_t,G_t).
\]

La politique peut produire des décisions discrètes, des paramètres
continus ou une combinaison des deux.

Le choix entre ces formulations dépend de la méthode d'apprentissage
et du type d'actions autorisées.

\subsection{Organisation interne du coordinateur}

\begin{figure}[H]
\centering
\begin{tikzpicture}[node distance=0.65cm]
  \node[data] (state) {État du système};

  \node[probe,right=1.0cm of state,yshift=1.0cm] (p1) {Sonde\\d'évaluation};
  \node[probe,right=1.0cm of state] (p2) {Sonde\\de sélection};
  \node[probe,right=1.0cm of state,yshift=-1.0cm] (p3) {Sonde\\d'organisation};

  \node[control,right=1.2cm of p2] (compose) {Composition\\des décisions};
  \node[data,right=1.0cm of compose] (action) {Actions};

  \draw[flow] (state) -- (p1);
  \draw[flow] (state) -- (p2);
  \draw[flow] (state) -- (p3);

  \draw[flow] (p1) -- (compose);
  \draw[flow] (p2) -- (compose);
  \draw[flow] (p3) -- (compose);
  \draw[flow] (compose) -- (action);

  \begin{scope}[on background layer]
    \node[neuralcore,fit=(p1)(p2)(p3)(compose),
      label=above:{\textbf{CORE NEURONAL DE COORDINATION}}] {};
  \end{scope}
\end{tikzpicture}
\caption{Le coordinateur est lui-même un core composé de sondes.}
\end{figure}

Les sondes du coordinateur peuvent être entraînées conjointement
ou partiellement séparées, selon l'architecture retenue.

\subsection{Le coordinateur apprend à modifier le graphe}

La caractéristique centrale d'EX9 est que le coordinateur peut
apprendre à modifier l'organisation du calcul.

Le graphe courant est noté :

\[
G_t=(V_t,E_t).
\]

Le coordinateur peut proposer une transformation :

\[
G_{t+1}
=
\mathcal{T}(G_t,a_t),
\]

où $\mathcal{T}$ est l'opérateur qui applique une action au graphe.

Les actions possibles comprennent :

\begin{itemize}
  \item \textbf{ACTIVER :} activer un noeud ou un core ;
  \item \textbf{DÉSACTIVER :} ne pas exécuter un noeud non nécessaire ;
  \item \textbf{AJOUTER :} ajouter une opération autorisée ;
  \item \textbf{RETIRER :} retirer une opération devenue inutile ;
  \item \textbf{RELIER :} ajouter une dépendance de données valide ;
  \item \textbf{DÉCONNECTER :} supprimer une dépendance devenue inutile ;
  \item \textbf{RÉORDONNER :} modifier l'ordre d'exécution ;
  \item \textbf{PARALLÉLISER :} lancer des branches indépendantes ;
  \item \textbf{RÉUTILISER :} employer un résultat déjà disponible ;
  \item \textbf{STOCKER :} conserver un résultat intermédiaire ;
  \item \textbf{LIBÉRER :} libérer une ressource ou une mémoire ;
  \item \textbf{ROUTER :} choisir un chemin de communication.
\end{itemize}

Une modification du graphe ne signifie pas que le coordinateur
peut créer arbitrairement n'importe quelle opération.

Les transformations doivent respecter les contrats fonctionnels
des cores et les dépendances du calcul.

\subsection{Graphe dynamique}

\begin{figure}[H]
\centering
\begin{tikzpicture}[node distance=0.7cm]
  % Graphe initial
  \node[probe] (a1) {Opération A};
  \node[probe,right=of a1] (b1) {Opération B};
  \node[probe,right=of b1] (c1) {Opération C};
  \node[data,right=of c1] (y1) {Sortie};

  \draw[flow] (a1) -- (b1);
  \draw[flow] (b1) -- (c1);
  \draw[flow] (c1) -- (y1);

  \node[below=0.35cm of b1] {\small Graphe initial};

  % Coordinateur
  \node[control,below=1.5cm of b1] (coord)
    {Coordinateur neuronal\\adaptatif};

  \draw[learned] (coord.north) -- (b1.south);

  % Graphe modifié
  \node[probe,below=1.6cm of a1] (a2) {Opération A};
  \node[probe,right=of a2] (b2) {Opération B};
  \node[probe,below=0.9cm of b2] (d2) {Opération D};
  \node[probe,right=of b2] (c2) {Opération C};
  \node[data,right=of c2] (y2) {Sortie};

  \draw[flow] (a2) -- (b2);
  \draw[flow] (b2) -- (c2);
  \draw[flow] (b2) -- (d2);
  \draw[flow] (d2) -| (c2);
  \draw[flow] (c2) -- (y2);

  \node[below=0.35cm of d2] {\small Graphe réorganisé};
\end{tikzpicture}
\caption{Illustration conceptuelle d'une adaptation de la structure
du graphe par le coordinateur.}
\end{figure}

Cette figure illustre une possibilité de réorganisation.
Elle ne définit pas une règle universelle de modification.

Une transformation est admissible uniquement si elle conserve
la validité fonctionnelle du calcul ou si elle correspond à une
transformation explicitement autorisée par la tâche.

% =========================================================
\section{Contraintes et validité de la coordination}
% =========================================================

\subsection{Pourquoi le graphe ne peut pas être modifié arbitrairement}

Un coordinateur neuronal peut apprendre à organiser le calcul,
mais ses décisions doivent respecter des contraintes.

Sans contraintes, il pourrait :
\begin{itemize}
  \item exécuter une opération avant que ses entrées existent ;
  \item supprimer une opération indispensable ;
  \item utiliser un résultat incompatible ;
  \item créer une dépendance circulaire ;
  \item dépasser les ressources disponibles ;
  \item produire une sortie de type incorrect.
\end{itemize}

Le système doit donc distinguer la capacité de proposer une action
et l'autorisation effective de l'exécuter.

\subsection{Ensemble d'actions admissibles}

On définit :

\[
\mathcal{A}_{\mathrm{valid}}(s_t,G_t)
\subseteq
\mathcal{A},
\]

où $\mathcal{A}$ est l'ensemble des actions théoriquement possibles.

Le coordinateur choisit parmi les actions admissibles :

\[
a_t
\in
\mathcal{A}_{\mathrm{valid}}(s_t,G_t).
\]

Une fonction de filtrage peut être définie :

\[
a_t
=
\operatorname{Mask}
\left(
\pi_{\theta_c}(s_t,G_t)
\right).
\]

Le filtrage doit notamment vérifier :
\begin{itemize}
  \item la compatibilité des types ;
  \item la disponibilité des entrées ;
  \item les dépendances nécessaires ;
  \item l'absence de cycle non autorisé ;
  \item les limites de ressources ;
  \item la validité des accès mémoire.
\end{itemize}

\subsection{Préservation de la sémantique}

Soit $G$ un graphe valide réalisant une fonction $F_G$.

Une transformation de graphe $\mathcal{T}$ est sémantiquement
équivalente si :

\[
F_{\mathcal{T}(G)}(x)
=
F_G(x)
\]

pour tout $x$ dans le domaine considéré.

Pour une transformation approximative, on peut exiger :

\[
d\left(
F_{\mathcal{T}(G)}(x),
F_G(x)
\right)
\leq \varepsilon,
\]

où $d$ est une distance adaptée à la sortie.

Cette condition est importante lorsque le coordinateur modifie
la structure du calcul.

Il faut distinguer :
\begin{itemize}
  \item les transformations exactes ;
  \item les transformations approximatives ;
  \item les transformations qui changent volontairement le calcul
        parce qu'elles correspondent à une décision apprise.
\end{itemize}

Ces trois catégories ne doivent pas être confondues dans les
expériences.

% =========================================================
\section{Apprentissage du coordinateur neuronal}
% =========================================================

\subsection{Objectif général}

Le coordinateur doit apprendre à produire des décisions qui
respectent trois objectifs distincts :

\begin{enumerate}
  \item \textbf{Qualité :} obtenir un résultat correct pour la tâche.
  \item \textbf{Efficacité :} limiter le temps, l'énergie, la mémoire
        et le travail exécuté.
  \item \textbf{Faisabilité :} respecter les dépendances, les contrats
        et les limites de ressources.
\end{enumerate}

La coordination n'est donc pas supposée être gratuite.
Elle constitue une fonction apprise dont le coût est lui-même
optimisé.

\subsection{Fonction objectif}

On définit une fonction objectif générale :

\[
\begin{aligned}
\mathcal{J}
={}&
\mathcal{L}_{\mathrm{task}}
+\lambda_T T_{\mathrm{total}}
+\lambda_E E_{\mathrm{total}}\\
&+\lambda_M M_{\mathrm{peak}}
+\lambda_R C_{\mathrm{redundancy}}
+\lambda_C C_{\mathrm{coord}}.
\end{aligned}
\]

Où :

\begin{itemize}
  \item $\mathcal{L}_{\mathrm{task}}$ est la perte de la tâche ;
  \item $T_{\mathrm{total}}$ est la latence ;
  \item $E_{\mathrm{total}}$ est l'énergie consommée ;
  \item $M_{\mathrm{peak}}$ est la mémoire maximale ;
  \item $C_{\mathrm{redundancy}}$ mesure le travail redondant ;
  \item $C_{\mathrm{coord}}$ mesure le coût propre de coordination.
\end{itemize}

Les coefficients $\lambda$ déterminent les compromis recherchés.

Ils doivent être fixés ou sélectionnés selon un protocole explicite.
Ils ne doivent pas être ajustés après observation des résultats
du jeu de test.

\subsection{Contraintes de faisabilité}

L'optimisation peut être formulée sous contraintes :

\[
\min_{\theta_c,\theta_p}
\mathbb{E}[\mathcal{L}_{\mathrm{task}}]
+
\lambda_T\mathbb{E}[T_{\mathrm{total}}]
+
\lambda_E\mathbb{E}[E_{\mathrm{total}}],
\]

sous les contraintes :

\[
\begin{aligned}
G_t &\in \mathcal{G}_{\mathrm{valid}},\\
M_{\mathrm{peak}} &\leq M_{\max},\\
T_{\mathrm{total}} &\leq T_{\max}
\quad \text{si une limite est imposée},\\
a_t &\in \mathcal{A}_{\mathrm{valid}}(s_t,G_t).
\end{aligned}
\]

Ici, $\theta_p$ représente les paramètres des cores de calcul
et $\theta_c$ ceux du coordinateur.

\subsection{Apprentissage conjoint}

Une possibilité est d'entraîner conjointement le coordinateur
et les cores :

\[
(\theta_c^\star,\theta_p^\star)
=
\arg\min_{\theta_c,\theta_p}
\mathcal{J}(\theta_c,\theta_p).
\]

Cette approche permet aux cores et au coordinateur de s'adapter
les uns aux autres.

Elle peut cependant être difficile à stabiliser, car une modification
du coordinateur change la distribution des opérations exécutées.

Une stratégie alternative consiste à :
\begin{enumerate}
  \item préentraîner les cores ;
  \item entraîner le coordinateur sur ces cores ;
  \item effectuer éventuellement un ajustement conjoint ;
  \item comparer cette stratégie à un entraînement conjoint direct.
\end{enumerate}

\subsection{Apprentissage par renforcement}

Les décisions de coordination peuvent être discrètes et avoir
des conséquences sur plusieurs étapes d'exécution.

On peut alors représenter le problème comme un processus
de décision séquentiel.

L'état est $s_t$, l'action est $a_t$, et la transition est :

\[
s_{t+1}
=
\mathcal{T}(s_t,a_t).
\]

Une récompense possible est :

\[
r_t
=
-\left(
\alpha c_{\mathrm{compute},t}
+\beta c_{\mathrm{memory},t}
+\gamma c_{\mathrm{communication},t}
+\delta c_{\mathrm{coord},t}
\right)
-r_{\mathrm{error},t}.
\]

La récompense finale peut inclure la qualité de la sortie.

L'objectif est de maximiser :

\[
J_{\mathrm{RL}}(\theta_c)
=
\mathbb{E}_{\pi_{\theta_c}}
\left[
\sum_{t=0}^{T}\gamma^t r_t
\right].
\]

L'apprentissage par renforcement est une option, pas une obligation.
Il devient particulièrement pertinent lorsque les décisions sont
discrètes, séquentielles et difficiles à différencier directement.

\subsection{Apprentissage supervisé à partir de traces}

Une autre stratégie consiste à produire des traces d'exécution
de référence.

Ces traces peuvent provenir :
\begin{itemize}
  \item d'un ordonnanceur conventionnel ;
  \item d'une recherche de graphe ;
  \item d'une optimisation hors ligne ;
  \item d'un planificateur exact sur de petits problèmes.
\end{itemize}

Le coordinateur apprend alors à imiter les décisions :

\[
\mathcal{L}_{\mathrm{coord}}
=
-\sum_t
\log
\pi_{\theta_c}
(a_t^\star\mid s_t,G_t).
\]

Cette approche permet de démarrer avec des comportements valides
avant d'optimiser directement les coûts.

\subsection{Différentiabilité}

Les paramètres continus des sondes peuvent être entraînés par
rétropropagation lorsque les opérations sont différentiables.

En revanche, certaines décisions du coordinateur sont discrètes :
\begin{itemize}
  \item activer ou désactiver un core ;
  \item ajouter ou retirer une arête ;
  \item choisir un ordre d'exécution ;
  \item sélectionner une opération.
\end{itemize}

Ces décisions ne sont pas directement différentiables dans leur
forme discrète.

Plusieurs méthodes sont envisageables :
\begin{itemize}
  \item apprentissage par renforcement ;
  \item apprentissage supervisé de décisions ;
  \item relaxation continue des choix ;
  \item estimateurs de gradient adaptés aux décisions discrètes.
\end{itemize}

Une relaxation différentiable peut faciliter l'entraînement,
mais elle doit être évaluée avec le comportement discret réellement
utilisé à l'inférence.

\subsection{Initialisation et démarrage}

Le coordinateur ne doit pas être supposé capable de construire
seul une architecture valide dès le début de l'apprentissage.

Une stratégie de démarrage peut être :

\begin{enumerate}
  \item définir un graphe initial valide ;
  \item initialiser les cores avec des contrats fonctionnels ;
  \item limiter les actions aux transformations sûres ;
  \item entraîner le coordinateur sur des exécutions de référence ;
  \item élargir progressivement son espace d'actions ;
  \item autoriser ensuite les modifications dynamiques du graphe.
\end{enumerate}

Cette stratégie évite de confondre l'apprentissage de la coordination
avec la découverte simultanée de toutes les fonctions élémentaires.

% =========================================================
\section{Langage d'actions du coordinateur}
% =========================================================

\subsection{Principe}

Le coordinateur manipule un langage d'actions décrivant les opérations
à effectuer sur les cores et le graphe.

On définit :

\[
\mathcal{A}
=
\{
\operatorname{APPLY},
\operatorname{ACTIVATE},
\operatorname{ROUTE},
\operatorname{STORE},
\operatorname{LOAD},
\operatorname{MERGE},
\operatorname{REWIRE},
\operatorname{WAIT},
\operatorname{RELEASE}
\}.
\]

Une action peut être représentée par un tuple :

\[
a_t =
(\mathrm{type},\mathrm{source},\mathrm{destination},
\mathrm{arguments}).
\]

L'exécuteur applique l'action si elle respecte les contrats
et les contraintes :

\[
\operatorname{Exec}(a_t,s_t)
=
s_{t+1}.
\]

Une séquence d'actions :

\[
A=(a_1,a_2,\ldots,a_T)
\]

produit une évolution du système :

\[
s_{t+1}
=
\operatorname{Exec}(a_t,s_t).
\]

Le résultat final est :

\[
y=H(s_T).
\]

\subsection{Exemple de programme d'actions}

Pour calculer :

\[
y=ab+\frac{e^c}{a}+\log(j),
\]

le coordinateur peut produire une séquence conceptuelle :

\[
\begin{aligned}
a_1 &= \operatorname{APPLY}(\mathrm{PROD},a,b),\\
a_2 &= \operatorname{APPLY}(\mathrm{EXP},c),\\
a_3 &= \operatorname{STORE}(u),\\
a_4 &= \operatorname{APPLY}(\mathrm{DIV},u,a),\\
a_5 &= \operatorname{APPLY}(\mathrm{LOG},j),\\
a_6 &= \operatorname{MERGE}(u_1,u_2,u_3).
\end{aligned}
\]

Le coordinateur peut aussi décider d'exécuter les opérations
indépendantes en parallèle.

Le langage d'actions ne supprime pas le calcul arithmétique.
Il décrit comment les opérations sont sélectionnées, ordonnées,
routées et réutilisées.

\subsection{Le programme d'actions est une sortie du coordinateur}

Il faut distinguer :

\[
\text{Coordinateur}
\neq
\text{Exécuteur}.
\]

Le coordinateur apprend à proposer les actions.
L'exécuteur applique les actions valides.

Cette séparation permet :
\begin{itemize}
  \item de vérifier les actions ;
  \item de mesurer le coût de coordination ;
  \item de comparer plusieurs politiques ;
  \item de remplacer le coordinateur sans réécrire les cores ;
  \item de garantir que les opérations exécutées respectent les contrats.
\end{itemize}

% =========================================================
\section{Architecture des cores spécialisés}
% =========================================================

\subsection{Core de calcul}

Le core de calcul regroupe les sondes qui réalisent les transformations
numériques ou neuronales nécessaires.

Il peut contenir :
\begin{itemize}
  \item des sondes arithmétiques ;
  \item des sondes RBF ;
  \item des projections ;
  \item des transformations vectorielles ;
  \item des opérations spécialisées.
\end{itemize}

\subsection{Core de composition}

Le core de composition organise les sorties de plusieurs sondes
ou cores.

Il peut réaliser :
\begin{itemize}
  \item des sommes ;
  \item des concaténations ;
  \item des projections ;
  \item des fusions conditionnelles ;
  \item des compositions hiérarchiques.
\end{itemize}

\subsection{Core de coordination}

Le core de coordination est un core neuronal adaptatif.

Il apprend à :
\begin{itemize}
  \item sélectionner les cores ;
  \item ordonner les tâches ;
  \item choisir les chemins d'exécution ;
  \item modifier le graphe ;
  \item gérer les ressources ;
  \item contrôler les réutilisations.
\end{itemize}

Il ne doit pas être confondu avec un simple ordonnanceur déterministe.

\subsection{Core de stockage temporaire}

Le core de stockage gère :
\begin{itemize}
  \item les résultats intermédiaires ;
  \item les états temporaires ;
  \item les clés de réutilisation ;
  \item les politiques de conservation ;
  \item l'invalidation des valeurs.
\end{itemize}

\subsection{Core de communication}

Le core de communication gère :
\begin{itemize}
  \item les transferts entre cores ;
  \item le routage des données ;
  \item la synchronisation ;
  \item la gestion des files ;
  \item les échanges entre unités physiques.
\end{itemize}

Il peut être intégré à d'autres cores lorsque son coût ne justifie
pas une séparation fonctionnelle.

\subsection{Organisation générale}

\begin{figure}[H]
\centering
\begin{tikzpicture}[node distance=0.8cm]
  \node[data] (input) {Entrée};

  \node[neuralcore,right=1.0cm of input] (coord) {
    \begin{minipage}{3.3cm}
      \centering
      \textbf{Coordinateur neuronal}\\
      Sondes adaptatives\\
      Politique apprise
    \end{minipage}
  };

  \node[core,below left=1.1cm and 1.0cm of coord] (calc) {
    \begin{minipage}{2.7cm}
      \centering
      \textbf{Core calcul}\\
      Opérations spécialisées
    \end{minipage}
  };

  \node[core,below=1.1cm of coord] (comp) {
    \begin{minipage}{2.7cm}
      \centering
      \textbf{Core composition}\\
      Fusion et assemblage
    \end{minipage}
  };

  \node[core,below right=1.1cm and 1.0cm of coord] (mem) {
    \begin{minipage}{2.7cm}
      \centering
      \textbf{Core mémoire}\\
      États et résultats
    \end{minipage}
  };

  \node[core,below=1.1cm of comp] (comm) {
    \begin{minipage}{2.7cm}
      \centering
      \textbf{Core communication}\\
      Routage et transfert
    \end{minipage}
  };

  \node[data,below=1.0cm of comm] (output) {Sortie};

  \draw[flow] (input) -- (coord);
  \draw[learned] (coord) -- (calc);
  \draw[learned] (coord) -- (comp);
  \draw[learned] (coord) -- (mem);
  \draw[learned] (coord) -- (comm);

  \draw[comm] (calc) -- (comp);
  \draw[comm] (comp) -- (mem);
  \draw[comm] (calc) -- (comm);
  \draw[comm] (mem) -- (comm);

  \draw[flow] (comm) -- (output);
\end{tikzpicture}
\caption{Organisation fonctionnelle proposée pour EX9.}
\end{figure}

Les flèches orange représentent les décisions de coordination.
Les flèches pointillées représentent les échanges entre cores.

Les cores présentés sont des rôles fonctionnels. Ils peuvent être
regroupés, distribués ou fusionnés selon le coût réel de leur
implémentation.

% =========================================================
\section{Architecture neuronale : remplacement des MLP}
% =========================================================

\subsection{MLP de référence}

Un bloc MLP classique peut être représenté par :

\[
\operatorname{MLP}(x)
=
W_2\sigma(W_1x+b_1)+b_2.
\]

Il effectue une transformation dense suivie d'une non-linéarité.

Dans EX9, un core neuronal doit apprendre à réaliser une fonction
comparable :

\[
C^{(N)}(x)
\approx
\operatorname{MLP}(x).
\]

\subsection{Core neuronal de remplacement}

Le core neuronal est composé de sondes adaptatives :

\[
C^{(N)}(x)
=
\mathcal{F}
\left(
P_1(x),P_2(x),\ldots,P_m(x)
\right).
\]

La fonction $\mathcal{F}$ peut être :
\begin{itemize}
  \item une somme pondérée ;
  \item une concaténation suivie d'une projection ;
  \item une composition hiérarchique ;
  \item un mécanisme de coordination appris ;
  \item une fusion conditionnelle.
\end{itemize}

Le coordinateur neuronal peut intervenir pour sélectionner
les sondes ou organiser leur composition.

Il faut cependant distinguer :
\begin{itemize}
  \item le core neuronal qui réalise la transformation ;
  \item le core de coordination qui décide comment l'organiser.
\end{itemize}

Ces deux cores peuvent être entraînés séparément ou conjointement.

\subsection{Comparaison fonctionnelle}

\begin{figure}[H]
\centering
\begin{tikzpicture}[node distance=0.65cm]
  \node[data] (x1) {$x$};
  \node[probe,right=of x1,minimum width=2.4cm] (l1) {Projection dense};
  \node[probe,right=of l1,minimum width=2.4cm] (act) {Non-linéarité};
  \node[probe,right=of act,minimum width=2.4cm] (l2) {Projection dense};
  \node[data,right=of l2] (y1) {$y$};

  \draw[flow] (x1) -- (l1);
  \draw[flow] (l1) -- (act);
  \draw[flow] (act) -- (l2);
  \draw[flow] (l2) -- (y1);

  \node[below=0.35cm of act] {\textbf{MLP de référence}};
\end{tikzpicture}

\vspace{0.8cm}

\begin{tikzpicture}[node distance=0.55cm]
  \node[data] (x) {$x$};
  \node[probe,right=of x,yshift=0.8cm] (p1) {Sonde 1};
  \node[probe,right=of x] (p2) {Sonde 2};
  \node[probe,right=of x,yshift=-0.8cm] (p3) {Sonde $m$};

  \node[control,right=1.1cm of p2] (f) {Composition\\adaptative};
  \node[data,right=of f] (y) {$y'$};

  \draw[flow] (x) -- (p1);
  \draw[flow] (x) -- (p2);
  \draw[flow] (x) -- (p3);
  \draw[flow] (p1) -- (f);
  \draw[flow] (p2) -- (f);
  \draw[flow] (p3) -- (f);
  \draw[flow] (f) -- (y);

  \node[below=0.35cm of f] {\textbf{Core neuronal EX9}};
\end{tikzpicture}
\caption{Remplacement expérimental d'un bloc MLP par un core neuronal.}
\end{figure}

\subsection{Protocole d'apprentissage}

Deux approches seront comparées.

\textbf{Apprentissage supervisé direct :}

\[
\Theta^\star
=
\arg\min_{\Theta}
\mathbb{E}_{(x,y)\sim\mathcal{D}}
\left[
\mathcal{L}(C_{\Theta}(x),y)
\right].
\]

\textbf{Distillation d'un MLP enseignant :}

\[
\Theta^\star
=
\arg\min_{\Theta}
\mathbb{E}_{x\sim\mathcal{D}}
\left[
D\left(
C_{\Theta}(x),
\operatorname{MLP}_{\mathrm{teacher}}(x)
\right)
\right].
\]

La distillation permet de comparer plus directement les fonctions
réalisées, indépendamment des différences d'optimisation.

% =========================================================
\section{Architecture globale et exécution EX9}
% =========================================================

\subsection{Définition du système}

Soit un ensemble de cores :

\[
\mathcal{C}
=
\{C_1,C_2,\ldots,C_M\}.
\]

On définit le système EX9 par :

\[
\operatorname{EX9}
=
\left(
\mathcal{C},
G_t,
C_{\mathrm{coord}},
\mathcal{M}_t,
\mathcal{N}
\right),
\]

où :
\begin{itemize}
  \item $\mathcal{C}$ est l'ensemble des cores ;
  \item $G_t$ est le graphe courant ;
  \item $C_{\mathrm{coord}}$ est le coordinateur neuronal ;
  \item $\mathcal{M}_t$ est l'état de la mémoire ;
  \item $\mathcal{N}$ représente les mécanismes de communication.
\end{itemize}

\subsection{Cycle général d'exécution}

Une inférence EX9 peut être organisée selon le cycle suivant :

\begin{enumerate}
  \item recevoir l'entrée ;
  \item construire ou mettre à jour sa représentation ;
  \item initialiser l'état du graphe ;
  \item transmettre l'état au coordinateur ;
  \item produire une décision de coordination ;
  \item vérifier la validité de l'action ;
  \item exécuter les opérations sélectionnées ;
  \item stocker ou réutiliser les résultats ;
  \item mettre à jour l'état du système ;
  \item répéter si d'autres actions sont nécessaires ;
  \item produire la sortie finale.
\end{enumerate}

\begin{figure}[H]
\centering
\begin{tikzpicture}[
  node distance=0.65cm,
  every node/.style={align=center}
]
  \node[data] (input) {Entrée};
  \node[control,below=of input] (state) {Construction de l'état};
  \node[control,below=of state] (coord) {Coordinateur neuronal};
  \node[control,below=of coord] (check) {Vérification des actions};
  \node[core,below=of check] (execute) {Exécution des cores};
  \node[memory,below=of execute] (memory) {Mémoire temporaire};
  \node[data,below=of memory] (output) {Sortie};

  \draw[flow] (input) -- (state);
  \draw[flow] (state) -- (coord);
  \draw[flow] (coord) -- (check);
  \draw[flow] (check) -- (execute);
  \draw[flow] (execute) -- (memory);
  \draw[flow] (memory) -- (output);

  \draw[comm] (memory.east) -- ++(1.4,0)
    |- (state.east);
  \draw[comm] (execute.east) -- ++(0.8,0)
    |- (coord.east);
\end{tikzpicture}
\caption{Cycle général d'une exécution EX9.}
\end{figure}

Le retour d'information vers le coordinateur permet de prendre
des décisions successives à partir de l'état actualisé.

Pour une tâche simple, une seule décision peut suffire.
Pour une tâche complexe, plusieurs étapes de coordination peuvent
être nécessaires.

\subsection{Pseudo-algorithme conceptuel}

\begin{verbatim}
ENTREE : x
ETAT   : s
GRAPHE : G

initialiser(s, G, x)

tant que non termine(s, G) :

    representation <- encoder_etat(s, G)

    action <- coordinateur_neuronal(representation)

    action_valide <- verifier_contraintes(action, s, G)

    s, G <- executer(action_valide, s, G)

    s <- mettre_a_jour_memoire(s)

sortie <- extraire_resultat(s)

retourner sortie
\end{verbatim}

Ce pseudo-algorithme décrit le principe général.
Il ne constitue pas une implémentation optimisée.

% =========================================================
\section{Modèle de coût}
% =========================================================

\subsection{Coût total d'exécution}

Le coût d'EX9 ne peut pas être évalué uniquement par le nombre
d'opérations arithmétiques.

On définit :

\[
\begin{aligned}
T_{\mathrm{total}}
={}&T_{\mathrm{compute}}
+T_{\mathrm{coord}}
+T_{\mathrm{compose}}\\
&+T_{\mathrm{memory}}
+T_{\mathrm{communication}}
+T_{\mathrm{validation}}.
\end{aligned}
\]

Le coût du coordinateur comprend notamment :
\begin{itemize}
  \item le calcul des sondes de coordination ;
  \item l'évaluation de l'état ;
  \item la production des actions ;
  \item la sélection et le filtrage des actions ;
  \item la mise à jour de la politique d'exécution.
\end{itemize}

Pour une réalisation matérielle, on ajoute :

\[
E_{\mathrm{total}}
=
E_{\mathrm{compute}}
+E_{\mathrm{coord}}
+E_{\mathrm{memory}}
+E_{\mathrm{communication}}.
\]

La réduction du calcul arithmétique peut être annulée par le coût
des transferts, de la coordination ou de la mémoire.

\subsection{Coût de coordination}

On définit :

\[
C_{\mathrm{coord}}
=
C_{\mathrm{inference}}
+
C_{\mathrm{decision}}
+
C_{\mathrm{validation}}
+
C_{\mathrm{update}}.
\]

Cette quantité mesure le travail nécessaire pour organiser
l'exécution.

Le coordinateur apprend à réduire ce coût, mais son apprentissage
ne le rend pas nul.

Une stratégie de coordination n'est avantageuse que si :

\[
T_{\mathrm{coord}}
+
T_{\mathrm{calcul\ organisé}}
<
T_{\mathrm{calcul\ de\ référence}}.
\]

\subsection{Réutilisation}

Soit $N_{\mathrm{naive}}$ le nombre d'exécutions d'opérations dans
une version sans réutilisation, et $N_{\mathrm{reuse}}$ le nombre
d'exécutions après partage des résultats.

On définit :

\[
R_{\mathrm{reuse}}
=
1-
\frac{N_{\mathrm{reuse}}}{N_{\mathrm{naive}}}.
\]

Cette mesure ne suffit pas à démontrer un gain de latence :
les opérations n'ont pas toutes le même coût.

Une mesure pondérée est :

\[
C_{\mathrm{ops}}
=
\sum_{i=1}^{N}c_i n_i,
\]

où $c_i$ est le coût estimé de l'opération $i$ et $n_i$ son
nombre d'exécutions.

\subsection{Coût de modification du graphe}

Une modification dynamique du graphe peut elle-même être coûteuse.

On définit :

\[
C_{\mathrm{graph}}
=
c_v N_{\mathrm{add}}
+
c_e N_{\mathrm{edge}}
+
c_r N_{\mathrm{remove}}
+
c_s N_{\mathrm{sync}},
\]

où :
\begin{itemize}
  \item $N_{\mathrm{add}}$ est le nombre de noeuds ajoutés ;
  \item $N_{\mathrm{edge}}$ est le nombre d'arêtes modifiées ;
  \item $N_{\mathrm{remove}}$ est le nombre de noeuds retirés ;
  \item $N_{\mathrm{sync}}$ est le nombre de synchronisations.
\end{itemize}

Les coefficients représentent les coûts associés à ces opérations.

Le coût de modification du graphe doit être intégré au coût
de coordination et mesuré séparément.

% =========================================================
\section{Hypothèses scientifiques falsifiables}
% =========================================================

Les hypothèses suivantes constituent des propositions expérimentales.
Elles ne sont pas considérées comme démontrées.

\begin{description}

  \item[H1 -- Réutilisation.]
  Le stockage des intermédiaires réduit le nombre de sous-calculs
  répétés sur les graphes comportant des dépendances partagées.

  \item[H2 -- Spécialisation.]
  Des sondes spécialisées peuvent réaliser certaines transformations
  avec un coût inférieur à une unité généraliste de référence.

  \item[H3 -- Composition.]
  Une composition explicite peut améliorer l'organisation et
  l'ordonnancement des opérations.

  \item[H4 -- Coordination adaptative.]
  Un coordinateur neuronal peut apprendre à sélectionner les cores
  et les opérations utiles en fonction de l'entrée et de l'état
  du système.

  \item[H5 -- Adaptation du graphe.]
  Un coordinateur neuronal peut apprendre des modifications de
  graphe qui améliorent le compromis entre qualité et coût
  d'exécution.

  \item[H6 -- Coordination économique.]
  L'optimisation conjointe de la politique de coordination et de
  son coût peut produire un système plus efficace qu'une politique
  statique ou qu'un coordinateur non optimisé.

  \item[H7 -- Core neuronal.]
  Un ensemble de sondes adaptatives peut remplacer certains blocs
  MLP à qualité comparable.

  \item[H8 -- Hiérarchie.]
  Une organisation hiérarchique de sondes, cores spécialisés et
  coordinateur neuronal peut être avantageuse sur certaines tâches.

  \item[H9 -- Réalisation physique.]
  Une correspondance physique entre cores et ressources matérielles
  peut améliorer le compromis latence-énergie sur des charges adaptées.

\end{description}

Chaque hypothèse doit être testée séparément.

% =========================================================
\section{Plan expérimental EX9}
% =========================================================

\subsection{Étape 1 : moteur de graphe logiciel}

Construire un moteur de graphe où :
\begin{itemize}
  \item chaque opération est un noeud ;
  \item chaque dépendance est une arête ;
  \item les résultats intermédiaires sont identifiables ;
  \item les opérations indépendantes peuvent être planifiées
        en parallèle ;
  \item les types et contrats des opérations sont vérifiés.
\end{itemize}

Comparer :
\begin{enumerate}
  \item calcul direct ;
  \item graphe sans réutilisation ;
  \item graphe avec réutilisation ;
  \item graphe avec exécution conditionnelle.
\end{enumerate}

\subsection{Étape 2 : cores logiques}

Créer les cores logiques suivants :
\begin{enumerate}
  \item core de calcul ;
  \item core de composition ;
  \item core de coordination ;
  \item core de stockage temporaire.
\end{enumerate}

Le core de communication peut être ajouté comme cinquième module
si son coût ou sa complexité justifie une séparation.

\subsection{Étape 3 : coordinateur neuronal}

Remplacer le coordinateur statique par un core neuronal composé
de sondes adaptatives.

Comparer au minimum :

\begin{enumerate}
  \item coordination statique ;
  \item coordination apprise sans modification du graphe ;
  \item coordination apprise avec sélection dynamique des cores ;
  \item coordination apprise avec modification du graphe.
\end{enumerate}

Cette comparaison permet de distinguer la contribution de la
sélection adaptative de celle de la modification de topologie.

\subsection{Étape 4 : apprentissage de la coordination}

Comparer plusieurs stratégies :
\begin{itemize}
  \item apprentissage supervisé sur des traces ;
  \item apprentissage par renforcement ;
  \item apprentissage conjoint avec les cores ;
  \item apprentissage en deux phases.
\end{itemize}

Mesurer :
\begin{itemize}
  \item la qualité des décisions ;
  \item la stabilité de l'apprentissage ;
  \item le nombre d'actions ;
  \item le coût de coordination ;
  \item la validité des graphes produits ;
  \item la généralisation à de nouvelles entrées.
\end{itemize}

\subsection{Étape 5 : core neuronal de remplacement}

Remplacer un bloc MLP par un core neuronal de sondes adaptatives.

Comparer :
\begin{itemize}
  \item MLP de référence ;
  \item core neuronal à sondes RBF ;
  \item core neuronal avec fusion pondérée ;
  \item core neuronal avec coordination adaptative.
\end{itemize}

Les budgets de paramètres, les données, les graines et les conditions
d'évaluation doivent être contrôlés.

\subsection{Étape 6 : coordination dynamique du graphe}

Autoriser progressivement le coordinateur à :
\begin{itemize}
  \item activer ou désactiver des opérations ;
  \item choisir entre plusieurs chemins valides ;
  \item modifier l'ordre d'exécution ;
  \item ajouter des opérations autorisées ;
  \item réutiliser des sous-calculs ;
  \item modifier les transferts entre cores.
\end{itemize}

Les modifications doivent être limitées à un ensemble d'opérations
valides et mesurables.

\subsection{Étape 7 : simulation physique}

Avant de concevoir une puce, simuler :
\begin{itemize}
  \item les unités de calcul ;
  \item les mémoires locales ;
  \item les transferts ;
  \item les délais de communication ;
  \item les ressources parallèles ;
  \item le coût du coordinateur.
\end{itemize}

Une simulation de coût matériel doit inclure les transferts et
la synchronisation, pas uniquement les opérations des sondes.

\subsection{Étape 8 : étude matérielle}

Si les résultats logiciels et simulés sont encourageants, étudier
une implémentation matérielle.

Les paramètres à mesurer comprennent :
\begin{itemize}
  \item surface logique ;
  \item fréquence ;
  \item latence ;
  \item énergie par opération ;
  \item bande passante mémoire ;
  \item coût de communication inter-cores ;
  \item coût de la coordination.
\end{itemize}

% =========================================================
\section{Protocole de comparaison}
% =========================================================

\subsection{Baselines}

Les comparaisons doivent inclure :

\begin{enumerate}
  \item MLP classique ;
  \item EX7 V003 ;
  \item graphe de calcul conventionnel ;
  \item graphe avec mémorisation ;
  \item EX9 à cores logiques ;
  \item EX9 avec core neuronal ;
  \item EX9 avec coordinateur neuronal ;
  \item EX9 avec coordinateur neuronal et graphe dynamique ;
  \item EX9 simulé comme architecture physique.
\end{enumerate}

\subsection{Mesures}

\begin{table}[H]
\centering
\small
\begin{tabularx}{\textwidth}{p{3.2cm}X}
\toprule
Mesure & Objectif \\
\midrule
Accuracy / F1 / log-loss
& Qualité prédictive selon la tâche \\

Nombre de paramètres
& Taille des modèles neuronaux \\

Latence p50, p95, p99
& Coût d'exécution réel \\

Nombre d'opérations
& Quantité de calcul effectuée \\

Taux de réutilisation
& Part des résultats intermédiaires réutilisés \\

Mémoire maximale
& Coût des états et intermédiaires \\

Coût de coordination
& Temps consommé par les décisions et l'ordonnancement \\

Nombre de modifications du graphe
& Intensité de l'adaptation structurelle \\

Taux d'actions invalides
& Fiabilité des décisions proposées \\

Énergie
& Coût énergétique par inférence, si mesurable \\

Robustesse
& Stabilité sur des entrées et distributions variées \\

Généralisation des politiques
& Capacité à coordonner des situations non rencontrées
durant l'entraînement \\
\bottomrule
\end{tabularx}
\caption{Mesures nécessaires à l'évaluation d'EX9.}
\end{table}

\subsection{Règles expérimentales}

Pour garantir une comparaison valide :
\begin{itemize}
  \item utiliser les mêmes données et partitions ;
  \item conserver les mêmes conditions matérielles ;
  \item rapporter plusieurs graines ;
  \item mesurer l'exécution complète ;
  \item inclure le coût de coordination ;
  \item inclure les transferts mémoire ;
  \item distinguer temps d'entraînement et temps d'inférence ;
  \item rapporter moyenne et dispersion des mesures ;
  \item séparer les données d'entraînement, validation et test ;
  \item rapporter les budgets de paramètres et de calcul.
\end{itemize}

La latence doit être mesurée avec un protocole identique pour
toutes les architectures, en distinguant si nécessaire le temps
de préparation, le temps de calcul et le temps total de bout en bout.

% =========================================================
\section{Ablations indispensables}
% =========================================================

Les ablations permettront d'identifier la contribution réelle
de chaque mécanisme.

\begin{longtable}{p{2.7cm}p{4.0cm}p{6.0cm}}
\toprule
Variante & Modification & Question testée \\
\midrule

EX9-A
& Sans mémoire temporaire
& La réutilisation réduit-elle le coût ? \\

EX9-B
& Sans coordination dynamique
& La sélection conditionnelle apporte-t-elle un gain ? \\

EX9-C
& Coordinateur à règles fixes
& L'apprentissage du coordinateur apporte-t-il un avantage ? \\

EX9-D
& Coordinateur neuronal sans modification du graphe
& La sélection adaptative suffit-elle ? \\

EX9-E
& Coordinateur neuronal avec graphe dynamique
& La modification de topologie apporte-t-elle un gain supplémentaire ? \\

EX9-F
& Sans composition hiérarchique
& La composition explicite est-elle utile ? \\

EX9-G
& Un seul core généraliste
& La spécialisation apporte-t-elle un avantage ? \\

EX9-H
& Core neuronal remplacé par MLP
& Le remplacement est-il compétitif ? \\

EX9-I
& Communication non optimisée
& Quel est le coût des transferts ? \\

EX9-J
& Coordination sans pénalité de coût
& L'optimisation explicite du coût de coordination est-elle utile ? \\

EX9-K
& Graphe statique
& L'adaptation structurelle est-elle utile ? \\

\bottomrule
\end{longtable}

% =========================================================
\section{Risques et limites}
% =========================================================

\subsection{Coût de coordination}

Un coordinateur neuronal peut lui-même être coûteux.

Il peut être plus lent qu'une exécution statique si ses décisions,
ses sondes et ses vérifications consomment plus de ressources
que les opérations qu'il permet d'éviter.

La coordination doit donc être optimisée et mesurée de bout en bout.

\subsection{Stabilité de l'apprentissage}

Un coordinateur qui modifie le graphe peut changer la distribution
des calculs réalisés.

Les cores peuvent alors recevoir des distributions d'entrées
différentes au cours de l'entraînement.

Cette interaction peut rendre l'apprentissage instable.

Il faudra étudier :
\begin{itemize}
  \item l'entraînement en plusieurs phases ;
  \item le gel temporaire de certains paramètres ;
  \item les contraintes sur les modifications de graphe ;
  \item les stratégies de mise à jour progressive.
\end{itemize}

\subsection{Validité des graphes}

Un graphe dynamique doit respecter les contrats fonctionnels
et les dépendances.

Une politique neuronale ne doit pas être autorisée à produire
des graphes arbitraires sans mécanisme de validation.

\subsection{Granularité}

Des sondes trop petites peuvent créer un excès de communication.
Des sondes trop grandes peuvent perdre l'avantage de spécialisation.

La granularité optimale doit être recherchée expérimentalement.

\subsection{Mémoire temporaire}

La mémorisation peut augmenter l'empreinte mémoire et introduire
des coûts de recherche, d'invalidation et de synchronisation.

\subsection{Approximation neuronale}

Un core neuronal peut produire une approximation différente du MLP.
Il faut mesurer la qualité, la stabilité et la généralisation.

\subsection{Difficulté du passage au matériel}

Une architecture logique efficace ne devient pas automatiquement
une architecture physique efficace.

Les contraintes de placement, de routage, de mémoire et d'énergie
peuvent modifier complètement le résultat.

% =========================================================
\section{Positionnement scientifique}
% =========================================================

EX9 combine plusieurs idées connues en informatique et en
apprentissage automatique :
\begin{itemize}
  \item graphes de calcul et graphes de dépendances ;
  \item unités de calcul spécialisées ;
  \item mémoire locale et réutilisation des résultats ;
  \item ordonnancement des tâches ;
  \item architectures parallèles ;
  \item modèles neuronaux modulaires ;
  \item apprentissage de politiques de routage ;
  \item modification dynamique de graphes.
\end{itemize}

La nouveauté éventuelle ne peut donc pas être déduite de la seule
présence de cores, de sondes ou d'un coordinateur neuronal.

Elle devra être établie par :
\begin{enumerate}
  \item une définition précise de l'organisation des sondes ;
  \item une définition précise du coordinateur neuronal ;
  \item une règle explicite de composition ou de modification
        des graphes ;
  \item une comparaison avec les architectures existantes ;
  \item des résultats reproductibles ;
  \item des ablations démontrant la contribution de chaque mécanisme.
\end{enumerate}

La contribution potentielle d'EX9 réside dans l'étude d'une
architecture unifiée dans laquelle :

\begin{itemize}
  \item les sondes constituent les éléments fonctionnels des cores ;
  \item les cores réalisent des fonctions spécialisées ;
  \item les graphes décrivent les dépendances et les possibilités
        de composition ;
  \item un core neuronal adaptatif apprend à organiser ces graphes ;
  \item la coordination peut modifier la structure du calcul
        en fonction des données et de l'état du système.
\end{itemize}

L'élément central à tester est donc la capacité d'un coordinateur
neuronal à apprendre non seulement quelles opérations exécuter,
mais également comment organiser leur coopération.

% =========================================================
\section{Feuille de route}
% =========================================================

\begin{table}[H]
\centering
\small
\begin{tabularx}{\textwidth}{p{1.8cm}p{4.2cm}X}
\toprule
Phase & Livrable & Critère de validation \\
\midrule

EX9.1
& Moteur de graphe de calcul
& Exactitude fonctionnelle et dépendances correctes \\

EX9.2
& Cores logiques
& Mesure des coûts de composition et d'exécution \\

EX9.3
& Mémoire temporaire
& Réutilisation correcte et gain mesuré \\

EX9.4
& Core neuronal
& Comparaison équitable avec MLP \\

EX9.5
& Coordinateur neuronal
& Apprentissage de décisions de coordination valides \\

EX9.6
& Graphe dynamique
& Modifications valides et gain mesuré \\

EX9.7
& Optimisation conjointe
& Compromis qualité, latence et mémoire \\

EX9.8
& Simulation physique
& Modèle de latence, mémoire et communication \\

EX9.9
& Étude matérielle
& Faisabilité et compromis énergie-performance \\

\bottomrule
\end{tabularx}
\caption{Plan de développement proposé pour EX9.}
\end{table}

% =========================================================
\section{Conclusion}
% =========================================================

EX9 propose de considérer le calcul comme un système organisé
d'opérations spécialisées et composables.

Les sondes constituent les éléments internes des cores.
Les cores sont spécialisés selon leur fonction :
calcul, composition, coordination, stockage ou traitement neuronal.

Trois formes de réalisation sont distinguées :
\begin{itemize}
  \item le core logique, qui organise les sondes dans un graphe logiciel ;
  \item le core physique, qui associe les sondes à des ressources matérielles ;
  \item le core neuronal, qui utilise des sondes adaptatives pour
        apprendre une transformation.
\end{itemize}

Le coordinateur constitue un élément central de cette architecture.
Il est lui-même un core neuronal composé de sondes adaptatives.

Il apprend à sélectionner les cores, organiser les opérations,
gérer les ressources, réutiliser les résultats et modifier le graphe
de calcul lorsque cela est utile et autorisé.

L'hypothèse centrale d'EX9 est qu'un système de cores spécialisés,
organisé par un coordinateur neuronal adaptatif, pourrait apprendre
des stratégies de calcul différentes selon les entrées et les
situations rencontrées.

Cette hypothèse implique deux questions expérimentales distinctes :

\begin{enumerate}
  \item Le coordinateur apprend-il à prendre de meilleures décisions
        que des stratégies statiques ou conventionnelles ?
  \item Le bénéfice de ces décisions dépasse-t-il le coût supplémentaire
        de la coordination, de la mémoire et des modifications du graphe ?
\end{enumerate}

La modularité et l'adaptation ne garantissent ni une accélération
ni une réduction énergétique. Les coûts de coordination, de mémoire,
de communication et de validation doivent être intégrés à chaque mesure.

EX9 doit donc être développé progressivement : d'abord sous forme
de graphe logiciel vérifiable, ensuite sous forme de cores logiques,
puis de cores neuronaux et d'un coordinateur adaptatif, enfin sous
forme de simulation et d'étude matérielles.

% =========================================================
\begin{thebibliography}{99}

\bibitem{ex7}
Projet EX,
\textit{Architecture V003 -- OuterCoord-Lite-7x1},
rapport technique interne, archive
\texttt{sondes\_ex7\_v003}.

\bibitem{rbf}
D. S. Broomhead and D. Lowe,
\textit{Multivariable Functional Interpolation and Adaptive Networks},
Complex Systems, vol. 2, pp. 321--355, 1988.

\bibitem{moe}
R. A. Jacobs, M. I. Jordan, S. J. Nowlan, and G. E. Hinton,
\textit{Adaptive Mixtures of Local Experts},
Neural Computation, vol. 3, no. 1, pp. 79--87, 1991.

\bibitem{flynn}
M. J. Flynn,
\textit{Some Computer Organizations and Their Effectiveness},
IEEE Transactions on Computers, vol. C-21, no. 9,
pp. 948--960, 1972.

\bibitem{dataflow}
J. B. Dennis,
\textit{First Version of a Data Flow Procedure Language},
Programming Symposium, Lecture Notes in Computer Science,
vol. 19, pp. 362--376, 1974.

\bibitem{roofline}
S. Williams, A. Waterman, and D. Patterson,
\textit{Roofline: An Insightful Visual Performance Model for
Multicore Architectures},
Communications of the ACM, vol. 52, no. 4, pp. 65--76, 2009.

\bibitem{rl}
R. S. Sutton and A. G. Barto,
\textit{Reinforcement Learning: An Introduction},
2nd edition, MIT Press, 2018.

\end{thebibliography}

\end{document}